\section{ClaudeにWhiteboards・Sheets・Data}
AnthropicのClaude公式アカウントは、共同で計画を整理するWhiteboards、スプレッドシートを作るSheets、旧Dashboardsを発展させたDataの展開開始を発表した。有料プラン向けのベータ提供を含み、ボードやトラッカー、データを使った説明資料の作成を支援するという。X上ではMicrosoftの業務ソフトとの競合を意識した反応も見られた。機能別の提供条件や企業向けの展開範囲は引き続き確認が必要だ。
\dailyxsource{Claude公式投稿}{https://x.com/claudeai/status/2120206365759094857}

\section{Fish Audio、双方向音声のLive 3}
Fish AudioのCEOは、聞きながら話し、相槌や割り込みにも対応するフルデュプレックス音声モデル「Live 3」を発表した。通常の音声合成にとどまらず、会話の途中で相手が発言した際の中断や確認、考え直しを扱えると説明している。初期の1,000ユースケースについて、カスタムエージェントを無料構築するとの案内もあった。実運用での遅延や既存音声モデルとの差は今後の評価対象となる。
\dailyxsource{Fish Audio関係者の発表}{https://x.com/rissa_cao/status/2120214987633467626}

\section{NVIDIA CEOを巡る疑惑と米議会対応}
NVIDIAのJensen Huang CEOについて、中国情報機関との接近疑惑を巡り米下院が説明を求め、応じない場合に召喚状を検討するとの投稿が拡散した。取締役会の独立調査や集団訴訟への言及もあったが、Grokが収集した範囲では議会の正式文書や一次報道への直接リンクは確認されていない。投稿上では召喚状は未発行で、違法行為も確認されていないとされる。疑惑の真偽や対中半導体政策への影響を確定的に扱うことは避けたい。
\dailyxsource{関連する二次投稿}{https://x.com/OssanJam/status/2120275158883815895}

\section{OpenAIの四半期売上報道が拡散}
The Informationの記事を紹介する投稿を通じ、OpenAIが投資家に2026年第3四半期の売上高を約100億ドル、前年同期比約50\%増と伝えたとの情報が広がった。企業向けのAI需要拡大と結び付けた議論が見られる一方、観測窓内でOpenAIの公式確認投稿は把握されていない。APIとサブスクリプションの内訳なども明らかではなく、金額と成長率は現段階では報道に基づく数値として扱う。
\dailyxsource{報道を紹介した投稿}{https://x.com/apples_jimmy/status/2120272795435454470}

\section{AnthropicのRSP閾値を巡る解釈}
AnthropicのIPO関連文書を引用し、AIによる研究開発の加速がResponsible Scaling Policy（RSP）の閾値に達したと判断された、との解説が投稿された。ただし、将来のモデルがもたらす加速を期待する評価であり、すでに同等の効果が実測されたという意味ではないとの留保が付いている。METRの報告との関連も論点となった。背景となるFable 5.2の正式発表は観測窓外であり、本号では窓内の解釈と反応を記録する。
\dailyxsource{RSP関連の解説投稿}{https://x.com/LeRoyDesCimes/status/2120269842934796481}

\section{METRを引用した超知能リスク論}
Redwood ResearchのBuck Shlegerisは、METRの研究などを参照し、1年以内に超知能が実現する可能性を約3分の1とする自身の見通しを長文で論じた。Anthropicの研究開発が約70\%加速したとの数値にも触れ、モデルの価値観が十分に整合しない場合の制御喪失リスクを懸念している。これらの確率や加速率は投稿者の評価・引用に基づくもので、観測窓内で独立に検証された測定結果ではない。
\dailyxsource{研究者の長文投稿}{https://x.com/bshlgrs/status/2120271165692461316}

\section{OpenAIのControl Roomに言及}
OpenAI公式アカウントは、研究インフラ上でモデルのワークロードを追跡し、安全性上の懸念が生じた際に迅速に停止する「Control Room」の構築について投稿した。出力を35秒以内に止められるとの説明がある。モデルの高性能化とともに、実行中の監視と停止能力を備える運用基盤への関心が高まっている。ただし、具体的な実装構成や適用範囲についての詳細資料はGrok収集内で確認できていない。
\dailyxsource{OpenAI公式投稿}{https://x.com/OpenAI/status/2119913172039868579}

\section{Muse Emberの費用対性能比較}
Metaの小規模モデル「Muse Ember」について、Claude Fable 5に近い性能を低コストで実現しているとの比較投稿が観測された。ベンチマークやデモ動画に触れる反応はあるものの、Grokはこの話題を未確認と評価している。公式のモデル仕様や測定条件、独立した比較結果が示されていないため、性能上の優劣を確定することはできない。小型モデルの実用性に対する関心を示す投稿として位置付けられる。
\dailyxsource{比較に関する投稿}{https://x.com/village_dev/status/2120274774731674016}
